\documentclass[11pt]{article}
\usepackage{amsmath, amssymb}
\usepackage{booktabs}
\usepackage{geometry}
\usepackage{hyperref}
\geometry{margin=1in}

\title{PPO Baseline on \texttt{markets-daily\_investor-v0}: Two-Seed Evaluation Report}
\author{Belief-Aware RL Project}
\date{May 2026}

\begin{document}
\maketitle

\section{Scope and motivation}

This note records the \emph{belief-blind} proximal policy optimization (PPO) baseline on
\texttt{markets-daily\_investor-v0}, trained with RLlib on the standard daily-investor observation
vector only (no belief features). The purpose is to anchor later comparisons against a reproducible
policy class $\pi_{\text{base}}(o_t)$, using the same reward definition (dense mark-to-market
increment, scaled in the environment) as in the design document \texttt{v3.tex}.

A third independent training seed was planned, but its run aborted during artifact I/O on the
shared filesystem; we therefore restrict this report to the \textbf{two completed seeds} (train
seeds $0$ and $1$). All conclusions are conditional on this two-sample slice of randomness.

\section{Training and evaluation protocol}

Each seed trains for $200{,}000$ environment timesteps with a $60$\,s trading-day timestep label,
$16$ parallel environment runners, train batch size $2000$, five PPO epochs per batch, minibatch
size $128$, learning rate $10^{-4}$, and entropy coefficient $0.01$. Checkpoints are written every
$50$ iterations. After training, the final policy is evaluated on \textbf{twenty} held-out
episodes drawn from evaluation seeds $100$--$119$ (inclusive), matching the plan-v1 harness.

Metrics are exported by \texttt{train\_ppo\_daily\_investor.py} into per-seed directories under
\texttt{results/ppo\_baseline/seed\_*}. Primary trading outcomes are episode-level final
profit-and-loss in \emph{cent} units as reported in \texttt{eval\_metrics.json} (denoted below as
final P\&L), together with Sharpe ratio, win rate, and mean maximum drawdown over episodes.

\section{Results}

\subsection{Aggregate view}

Both seeds complete the full training budget and pass the held-out evaluation suite without
degenerate failure modes. Mean final P\&L is strictly positive in each seed, with comparable
dispersion across the twenty evaluation episodes. Win rate is $1.0$ in both cases under the
exported definition (every evaluated episode records non-negative final P\&L in this run).

Across the two seeds, mean final P\&L averages $32{,}853.4$ cents per episode, with seed $1$
exceeding seed $0$ by about $4{,}651$ cents on average. Sharpe ratios are strong in both runs but
\emph{not} identical: seed $0$ attains the higher risk-adjusted score, while seed $1$ attains the
higher mean terminal P\&L, illustrating the usual mean--volatility tradeoff at small $n_{\text{seed}}$.

\subsection{Per-seed metrics}

Table~\ref{tab:two-seed} summarizes training wall time, evaluation wall time, and the main
evaluation statistics for each completed seed.

\begin{table}[h]
\centering
\begin{tabular}{lrr}
\toprule
\textbf{Quantity} & \textbf{Seed $0$} & \textbf{Seed $1$} \\
\midrule
Train wall time (s) & $1597.8$ & $1554.4$ \\
Eval wall time (s) & $654.2$ & $676.0$ \\
Mean final P\&L (cents) & $30{,}528.1$ & $35{,}178.7$ \\
Std.\ final P\&L (cents) & $11{,}352.9$ & $15{,}134.8$ \\
Min / max final P\&L (cents) & $10{,}941$ / $56{,}394$ & $6{,}946$ / $59{,}655$ \\
Mean episode total reward & $7.83$ & $9.02$ \\
Sharpe (episode level) & $2.689$ & $2.324$ \\
Win rate & $1.00$ & $1.00$ \\
Mean max drawdown & $0.00298$ & $0.00341$ \\
\bottomrule
\end{tabular}
\caption{Two-seed PPO baseline: training cost and held-out evaluation on twenty episodes per seed.}
\label{tab:two-seed}
\end{table}

\subsection{Action mix}

The learned policies allocate roughly half of decisions to aggressive sides (buy and sell) with
the remainder in hold, but the \emph{hold} share differs materially between seeds: seed $0$ uses
hold about $4.9\%$ of evaluated actions, whereas seed $1$ uses hold about $12.0\%$. Buy and sell
masses remain near balance in both cases, which is consistent with a market-making style that
leans on round-trip inventory risk rather than a persistent directional bet.

\subsection{Interpretation}

Taken together, the two seeds support a conservative headline: the PPO baseline under the stated
hyperparameters reliably discovers policies that are profitable on the held-out evaluation slice
and exhibit modest drawdowns relative to terminal P\&L. We do \emph{not} claim convergence to a
unique policy: action statistics and Sharpe versus mean P\&L differ nontrivially between seeds.

Because $n_{\text{seed}} = 2$, any aggregate (mean-of-seed-means) has large uncertainty; the third
seed should be resumed once stable scratch space is available, and \texttt{summary.json} should be
regenerated from three seeds for plan-v1 parity.

\section{Supplement: targeted heuristic baseline sweep}

This section summarizes a separate \emph{belief-free} diagnostic run: a one-at-a-time sweep of four
\texttt{rmsc04} background knobs (\texttt{num\_value\_agents}, \texttt{num\_momentum\_agents},
\texttt{kappa\_oracle}, \texttt{fund\_vol}), each at three values, crossed with four fixed heuristic
trading policies (\textsc{Noise}, \textsc{Value}, \textsc{Momentum}, \textsc{MarketMaker}) implemented in
\texttt{abides-gym/scripts/sweep.py}. Each $(\text{parameter}, \text{value}, \text{strategy})$ cell uses
ten episodes (the script's \texttt{--fast} mode) with a $300$\,s simulator timestep label. All rows are
aggregated in \texttt{targeted\_sweep.json} at the project root ($480$ rows total: $12$ cells $\times$
four strategies $\times$ ten seeds).

\paragraph{Robustness and timeouts.}
Raw mean P\&L can be dominated by a handful of microstructure blow-ups (especially with many momentum
followers). The discussion below therefore emphasizes \emph{medians} over the ten episode seeds within each
cell, and excludes the eight episodes that hit a hard wall-clock cap (all eight occurred for
\texttt{num\_momentum\_agents}${}=50$, under the episode-isolation guard added after an observed hang in
exchange order-book imbalance computation). Those capped episodes are retained in the JSON with
\texttt{timed\_out=true} and are omitted from the medians here.

\paragraph{Regime sensitivity (\texttt{num\_momentum\_agents}).}
Table~\ref{tab:sweep-mom} contrasts the three momentum-population settings while other background fields
follow the sweep's per-cell overrides. Raising momentum participation from $0$ to $12$ leaves the
\textsc{Value} / \textsc{Momentum} median split similar to the calibrated default ($12$ agents), whereas
$50$ momentum agents produces sharply negative median outcomes for \textsc{Momentum} and
\textsc{MarketMaker} and heavy left tails for \textsc{Noise} / \textsc{Value}---consistent with a crowded
feedback-trading regime that is hostile to simple heuristics.

\begin{table}[h]
\centering
\begin{tabular}{lrrrr}
\toprule
\texttt{num\_momentum\_agents} &
\multicolumn{2}{c}{\textsc{Value}} &
\multicolumn{2}{c}{\textsc{Momentum}} \\
\cmidrule(lr){2-3} \cmidrule(lr){4-5}
& median P\&L & win \% & median P\&L & win \% \\
\midrule
$0$   & $328$   & $50$ & $-1967$ & $30$ \\
$12$  & $-2485$ & $30$ & $1690$  & $60$ \\
$50$  & $-344{,}390$ & $50$ & $-1{,}561{,}646$ & $0$ \\
\bottomrule
\end{tabular}
\caption{Median final P\&L (cents) and episode win rate within each \texttt{num\_momentum\_agents} cell,
\textsc{Value} vs.\ \textsc{Momentum}, excluding timed-out episodes (typically $n=10$ episodes per
strategy; $n=7$ for \textsc{Momentum} at fifty agents after removing three timed-out runs).}
\label{tab:sweep-mom}
\end{table}

\paragraph{Oracle tightness (\texttt{kappa\_oracle}).}
Across the three oracle-gain settings, median \textsc{Value} final P\&L stays mildly negative (roughly
$-2.4$k to $-2.1$k cents), while \textsc{Momentum} medians range from about $-1.1$k to $+2.5$k cents. The
middle setting $1.67\times 10^{-16}$ matches the default \texttt{rmsc04} calibration and overlaps the
\texttt{num\_momentum\_agents}${}=12$ row in Table~\ref{tab:sweep-mom} because that sweep block reuses the
same background template.

\paragraph{Fundamental volatility (\texttt{fund\_vol}).}
Median outcomes move with volatility: the highest tested \texttt{fund\_vol} ($5\times 10^{-4}$) raises
\textsc{Momentum}'s median to about $+4.4\times 10^{4}$ cents (eight of ten seeds profitable) while pushing
\textsc{Value} and \textsc{MarketMaker} medians further negative, suggesting that ``buy the drift'' heuristics
benefit more than mean-reversion-style rules when shocks are large under this microstructure.

\paragraph{Takeaway for the RL agenda.}
The sweep is not a substitute for trained policies, but it quantifies \emph{where} the daily-investor MDP
becomes harsh or unstable for fixed baselines---especially under many momentum agents---and motivates
belief-aware modeling (Section~2 of \texttt{v3.tex}) as a way to condition on latent composition rather than
assuming a single benign background.

\end{document}
